\section{序論}
\label{sec:introduction}
経済活動や需要の多変量時系列では、他の情報を条件とした予測上の必要性から系列間の関係を調べるグレンジャー因果分析が用いられる\citep{granger1969}。線形ベクトル自己回帰（VAR）は係数から関係を読めるが、非線形モデルでは入力が予測に関与する複数の経路を制御する必要がある。

この課題に対し、\citet{tank2022}は入力層の構造正則化によって依存構造を抽出した。一方、一般化VAR（generalised VAR; GVAR）は、過去値を入力依存の係数で重み付けし、非線形な関係を係数として表す\citep{marcinkevics2021}。この係数表現で不要な入力を除くには、過去値に掛かる係数だけでなく、その入力が他の係数を変える経路も制御する必要がある。

そこでGated Self-Explaining VAR（GSE-VAR）を提案する。対象・入力・ラグ別の共通ゲートで係数生成の入力と出力を制御し、各ゲートの選択確率に期待L0正則化を課し、DisARMで期待損失のロジット勾配を推定する。係数生成器とロジットはAdamで同時に更新する。

本研究の貢献は、非線形な係数表現と、ゲートの\(0\)による入力非依存性の判定を両立する設計にある。十分条件と符号解釈を示し、Lorenz96の構造回復と線形VARの符号回復から評価する。さらに付録\ref{sec:application}では、Favoritaの販売系列を用いて、正解構造が未知の場合のモデル選択とカテゴリ間関係の記述を例示する。
